\section{The actual function-field degree of a finite dominant map}

\begin{definition}[The canonical original generic extension degree]
\label{def:intrinsic-numerical-canonical-function-field-degree}
\uses{lem:actual-dominant-function-field-map}
\lean{EnvelopeProperty.NumericalCurveFieldDegree.degree}
\leanok
For an actual finite dominant $g:D\to C$ of integral schemes, let
$N(g)=[k(D):k(C)]$, computed with the actual generic-stalk field
embedding induced by $g$. This definition uses the whole original
function fields and that same canonical embedding.
\end{definition}

\begin{theorem}[Actual finiteness makes this original degree positive]
\label{thm:intrinsic-numerical-function-field-degree-positive}
\uses{def:intrinsic-numerical-canonical-function-field-degree,lem:actual-dominant-stalk-injectivity,thm:fc-scheme-finite-principal-generic-rank}
\lean{EnvelopeProperty.NumericalCurveFieldDegree.functionFields_finite,EnvelopeProperty.NumericalCurveFieldDegree.degree_pos}
\leanok
If the same $g$ is finite, then its actual function-field extension is
finite and $N(g)>0$. No flatness, normality, dimension bound or
independently chosen field embedding is required.
\end{theorem}
\begin{proof}
Choose a nonempty original affine open $U$ of $C$. Dominance makes
$g^{-1}U$ nonempty, and finiteness makes it affine. Write
$R=\Gamma(C,U)$ and $B=\Gamma(D,g^{-1}U)$. These are original domains,
and the actual section map makes $B$ finite, hence integral, over $R$.
The original fields $K=k(C)$ and $F=k(D)$ are respectively the fraction
localizations of $R$ and $B$. Naturality of the actual generic germs
identifies the two maps from $R$ to $F$: the map through $B$ equals the
map through the canonical original $K\to F$.

Integrality of $B/R$ and fraction localization make $F$ algebraic over
$K$. The finite-type algebra $B/R$, followed by its fraction
localization, makes $F$ essentially of finite type over $R$, and then
over $K$. An algebraic field extension essentially of finite type is
finite: finitely many field generators are algebraic, so their field
adjoin equals their finite algebra adjoin. Thus the actual $F/K$ is
finite. Its underlying vector space is nonzero, so its dimension,
which is $N(g)$, is positive. No assertion that $K$ is a finite-type
$R$-algebra is used.
\end{proof}

\begin{theorem}[Every original affine inverse computes the same degree]
\label{thm:intrinsic-numerical-function-field-degree-affine-rank}
\uses{def:intrinsic-numerical-canonical-function-field-degree,lem:actual-dominant-stalk-injectivity,thm:fc-scheme-finite-principal-generic-rank}
\lean{EnvelopeProperty.NumericalCurveFieldDegree.degree_eq_affineRank}
\leanok
For every nonempty original affine open $U$ of $C$ and the same actual
finite dominant $g$,
\[
 N(g)=\operatorname{rank}_{\Gamma(C,U)}\Gamma(D,g^{-1}U),
\]
with the scalar action given by the actual original section map.
\end{theorem}
\begin{proof}
Use the same actual rings and fraction fields as above. The coefficient
identity for the actual generic germs supplies both original scalar
towers into $F$. Fraction localization preserves module rank in each
tower, giving
$[F:K]=\operatorname{rank}_R B$. Both sides are computed using these
actual original maps. The argument works for each original nonempty
affine open separately and needs no normality or flatness.
\end{proof}
